\section{Mentalidad de investigación: negar cada día a quien uno fue ayer}

Luo Fuli (罗福莉) dice que lo que más ha sentido en el último medio año es que «cada día niega a quien fue ayer». No es una expresión emocional, sino una mentalidad de investigación ante un cambio de paradigma acelerado: cuando OpenClaw, Claude Code, el posentrenamiento de Agent y la orquestación de múltiples modelos aparecen con rapidez, los juicios anteriores pierden validez en cuestión de semanas o incluso de días. El equipo debe ser capaz de actualizar hipótesis, rehacer experimentos y reajustar la asignación de recursos con rapidez.

\subsection{El cambio de reward: de la inversión cuantitativa a los modelos grandes}

En la entrevista se menciona que, en la inversión cuantitativa, el precio es el reward, que resulta relativamente claro; en cambio, en el campo de los modelos grandes el reward es menos claro y más cambiante. Con los Agent esto se acentúa: el éxito de la tarea, el valor para el usuario, los límites de seguridad, la eficiencia en costos y el aprendizaje a largo plazo pueden formar parte del reward. Esto exige que el investigador no se limite a optimizar un solo número, sino que defina de manera continua «qué es lo bueno».

\begin{knowledgebox}{Cómo investigar cuando el reward no es claro}
Cuando el reward no es claro, la investigación no puede detenerse; primero hay que construir indicadores sustitutos verificables: tasa de finalización de tareas, preferencia humana, costo, tipos de falla, capacidad de recuperación, incidentes de seguridad, retención de usuarios, entre otros. Después hay que comprobar de forma constante si esos indicadores sustitutos realmente sirven al valor final.
\end{knowledgebox}

\subsection{El entorno importa más que la experiencia}

«El entorno importa más que la experiencia» es el juicio estratégico con el que cierra este episodio. Aquí el entorno incluye la capacidad de cómputo, la infraestructura, los chips, los sistemas operativos, el tráfico, la dimensión social, la cultura organizacional y los recursos estratégicos. La experiencia es útil, pero si el entorno no da retroalimentación, no permite equivocarse y probar, no proporciona capacidad de cómputo ni apoya la colaboración entre equipos, la experiencia se convierte pronto en un lastre del paradigma anterior.

\begin{importantbox}{El entorno es la aceleración de la investigación}
La experiencia determina la posición inicial; el entorno determina la aceleración. En la era de los Agent, los cambios son demasiado rápidos: quien obtenga antes retroalimentación real, más rollout, mejor infra y una organización más abierta tendrá más probabilidades de liderar en el nuevo paradigma.
\end{importantbox}

\subsection{Resumen de la sección}

Negar cada día a quien uno fue ayer no significa carecer de juicio, sino construir un sistema capaz de actualizar sus juicios. La mentalidad de investigación en la era de los Agent consiste en enlazar valores, retroalimentación experimental, entorno de recursos y ajustes organizacionales en un ciclo continuo de aprendizaje.
